\documentclass[10pt]{article}

\usepackage[utf8]{inputenc}

\usepackage[T1]{fontenc}

\usepackage{amsmath}

\usepackage{amsfonts}

\usepackage{amssymb}

\usepackage[version=4]{mhchem}

\usepackage{stmaryrd}

\usepackage{bbold}



\begin{document}

го квантования может быть получен из вакуумного действием на него оператора рождения частиц. В ряде случаев, напр. при спонтанном нарушении симметрии, вакуумное состояние оказывается не единственным, вырожденным,существует непрерывный спектр таких состояний, отличающихся друг от друга числом так наз. голдстоуновских бозонов. А. В. Ефремов. ВАКУУМНАЯ ДИАГРАММА ФЕЙНМАНА - сМ. ФеЙнмана диаграмма.



ВАКУУМНОЕ КЛАСТЕРНОЕ РАЗЛОЖЕНИЕ - СМ. КЛастерное разложение.



BAKУУMНОЕ СОСТОЯНИЕ - см. Вакуум.



BAKУУМНОЕ СРЕДНЕЕ в к в а н то в о Й т е о и и пол я - комплексное число, равное среднему значению какого-либо оператора (или произведения операторов $A, B, \ldots$ ) в вакуумном состоянии поля $|0\rangle$ (см. Вакуум). Обозначается символом $\langle 0|A ; B ; \ldots| 0\rangle$. В. с. операторов энергии, импульса, момента импульса, электрич. заряда и др. сохраняюшихся квантовых чисел равны нулю. Особенно большую роль играет В. с. локальных операторов поля $\varphi(x)$, зависящих от пространственно-временных точек $x$. Так, ненулевое значение $\langle 0|\varphi(x)| 0\rangle$ свидетельствует о спонтанном нарушении симметрии и вырождении вакуума. В.с. от хронологич. произведения операторов полей или локальных токов дает матричные элементы матрицы рассеяния и определяет все процессы взаимопревращения частиц. ВАКУУМНЫЙ КОНДЕНСАТ - см. Квантовая хромодинамика; правила сумм.



ВАЛЕНТНОСТь те нзо ра - см. Тензорное поле.



ВАН ДЕР ПОЛЯ УРАВНЕНИЕ - нелинейное обыкновенное дифференциальное уравнение 2-го порядка:



\[

\ddot{x}-\mu\left(1-x^{2}\right) \dot{x}+x=0, \quad \mu=\text { const }>0, \quad \dot{x}(t) \equiv \frac{d x}{d t} .

\]



Является важным частным случаем Льенара уравнения. В. д. П.у. описывает свободные автоколебания одной из простейших нелинейных колебательных систем (ос ци ллятора Ван де р Поля). В частности, уравнение (1) служит математич. моделью (при ряде упрощающих предположений) лампового генератора на триоде в случае кубич. характеристики лампы. Характер решений уравнения (I) был впервые подробно изучен Б. Ван дер Полем (B. Van der Pol, 1922-26).



Уравнение (1) эквивалентно системе двух уравнений относительно фазовых переменных $x, v$ :



\[

\dot{x}=v, \quad \dot{v}=-x+\mu\left(1-x^{2}\right) v .

\]



Иногда вместо $x$ удобнее ввести переменную



\[

z(t)=\int_{0}^{t} x(\tau) d \tau

\]



тогда уравнение (1) приведется к уравнению



\[

\ddot{z}-\mu\left(\dot{z}-\frac{\dot{z}^{3}}{3}\right)+z=0

\]



являющемуся частным случаем Рэлея уравнения. Если вместе с переменной $x$ рассмотреть переменную



\[

y=-x+\frac{x^{3}}{3}+\frac{\dot{x}}{\mu}

\]



ввести новое время $\tau=t / \mu$ и положить $\varepsilon=\mu^{-2}$, то вместо уравнения (1) получим систему



\[

\varepsilon x^{\prime}=y-x+\frac{x^{3}}{3}, \quad y^{\prime}=-x, \quad '=\frac{d}{d t} .

\]



При любом $\mu>0$ в фазовой плоскости системы (2) сушествует единственный устойчивый предельный цикл, к к-рому при $t \rightarrow \infty$ приближаются все остальные траектории (кроме положения равновесия в начале координат); этот предельный цикл адекватен автоколебаниям осциллятора Ван дер Поля (см. [1]-[3]).



\section{BAKYУMTAR}

При малых $\mu$ автоколебания осциллятора (1) близки к простым гармонич. колебаниям с периодом $2 \pi$ и с определенной амплитудой. Для вычисления колебательного процесса с большей точностью применяются асимптотич. методы. При возрастании $\mu$ автоколебания осциллятора (1) все более отклоняются от гармонич. колебаний. При больших $\mu$ уравнение (1) описывает релаксационные колебания с периодом (в первом приближении) 1,614 $\mu$. Известны более точные асимптотич. разложения величин, характеризующих релаксационные колебания (см. [4]); изучение этих колебаний равносильно исследованию решений системы (3) с малым параметром $\varepsilon$ при производной (см. [5]).



Уравнение



\[

\ddot{x}-\mu\left(1-x^{2}\right) \dot{x}+x=E_{0}+E \sin \omega t

\]



описывает поведение осциллятора Ван дер Поля под воздействием внешнего периодич. возмущения. Здесь наиболее важны изучение явления зах ваты в ан и я ч а с тоты (существования периодич. колебаний) и исследование б и е н и й (возможности почти периодич. колебаний; см. [1], [3]).



Лum.: [1] Андрон о в А. А., В итт А. А., Хайкин С. Э., Теория колебаний, 2 изд., М., 1959; [2] Л е фш е ц С., Геометрическая теория дифференциальных уравнений, пер. с англ., М., 1961; [3] С ток е р Дж., Нелинейные колебания в механических и электрических системах, пер. с англ., 2 изд., М., 1953; [4] Дородн ицын А. А., «Прикл. математика и механика», 1947, т. 11, с. 313-28; [5] М ище н ко Е. Ф., Ро з о в Н.Х., Дифференциальные уравнения с малым параметром и релаксационные колебания, М., $1975 . \quad$ Н.Х. Розов. BAH XOBA TEOPEMA - теорема, устанавливающая отсутствие фазовых переходов в одномерных непрерывных статистических системах частиц с парным потенциалом взаимодействия.



Пусть парный потенциал взаимодействия $\Phi: \mathbb{R} \rightarrow \mathbb{R} \cup\{+\infty\}$



a) непрерывен при $|x|>R$;



б) имеет твердую сердцевину $[\Phi(x)=+\infty$ при $|x|<R$, $x \in \mathbb{R}]$



в) имеет конечный радиус действия (финитен) $\Phi(x) \equiv 0$ при $|x| \geqslant R_{0}>R$.



В этом случае предельное $(V(\Lambda) \rightarrow \infty, V(\Lambda)$ - объем области $\Lambda$, в к-рой находится система) давление



\[

p(z, \beta)=\beta^{-1} \lim _{\Lambda \rightarrow \infty} V(\Lambda)^{-1} \ln \Xi(\Lambda, z, \beta),

\]



где $\Xi-$ большая статистич. сумма, является действительной аналитич. функцией активности $z$ (при $z>0$ ). Следовательно, в таких системах фазовые переходы отсутствуют.



Впоследствии класс потенциалов взаимодействия, для к-рых справедлива В.Х. т., был существенно расширен. Так, в частности, В. Х. т. справедлива (см. [3], [5]) для быстро убывающих потенциалов взаимодействия с твердой сердцевиной, удовлетворяющих условию $|\Phi(x)| \leqslant \varphi(|x|)$ при $|x|<A$, где $\varphi(r)$ - монотонно невозрастающая функция такая, что



\[

\int_{A}^{\infty} r \varphi(r) d r<\infty

\]



а также (см. [6]) для финитных потенциалов взаимодействия, быстро растущих в нуле.



Существует также В.Х. т., устанавливающая наличие и характер особых точек в энергетич. зонах кристаллов.



Jum.: [1] Van Hove L., "Physica», 1950, t. 16, p. 137-43; [2] Р ю эл ь Д., Статистическая механика. Строгие результаты, пер. с англ., М., 1971, теоремы 5, 6, 7 и § 4 прилож.; [3] Доб руш и н Р. Л., «Функциональный анализ и его прилож.», 1969, № 1, с. 27; [4] е го же, «Теория вероятн. и ее примен.», 1970, т. 3, с. 469; [5] Gallavotti G., Miracle-Sole S., «J. Math. Phys.», 1970, v. 1, р. 147; [6] Сухов Ю. М., «Успехи математич. наук», 1970, т. 25, B. 2 , c. 277.



П. В. Малышев.



BAPAДАНА TEOPEMA - теорема, определяющая асимптотическое поведение интегралов по последовательности вероятностных мер, удовлетворяюших больших уклонений принципу. Фактически она обобщает метод Лапласа на бесконечномерный случай. Наиболее удобным для различных приложений в физике является следуюший вариант В.т.





\end{document}